\documentclass[conference]{IEEEtran}
\usepackage{cite,amsmath,amssymb,booktabs,array}
\newcommand{\indicator}{\mathbb{I}}
\begin{document}
\title{Can Humans Hibernate? A Staged Proposal for a Human Torpor Safety Envelope}
\author{\IEEEauthorblockN{Anonymous Research Plan Author}\\\IEEEauthorblockA{Frozen Survey--Idea--ExperimentDesign record\\Design-only scientific research plan}}
\maketitle
\begin{abstract}Human hibernation is not established by lowering body temperature. Natural torpor is an actively controlled, reversible organism-level state, whereas uncontrolled hypothermia can cause arrhythmia, impaired perfusion, coagulopathy, metabolic instability, infection risk, and injury during rewarming. This proposal defines a falsifiable research program for a bounded human torpor safety envelope. It compares closed-loop thermometabolic control with matched normothermic protocols, measures metabolic expenditure and safety continuously, and forbids escalation after any preregistered failure. The work is design-only: no human hibernation, efficacy result, or clinical recommendation is claimed.\end{abstract}
\begin{IEEEkeywords}hibernation, torpor, hypothermia, thermoregulation, metabolism, closed-loop control, rewarming\end{IEEEkeywords}
\section{Introduction}
The claim that cooling could create a state requiring no fuel is biologically false. Metabolic demand may decline as temperature declines, but ion gradients, circulation, ventilation, immune defense, and cellular repair continue to require energy. Hibernators coordinate these functions through evolved programs; a cooled human does not automatically enter one. The scientific question is therefore whether a reversible, low-metabolic state can be approached within explicit safety constraints.
Natural torpor and hibernation involve reductions in metabolic rate and temperature that differ across species and seasons \cite{geiser2004,ruf2015}. Human targeted temperature management is useful as a clinical comparison but not a proof of hibernation. Hypothermia has recognized systemic complications and rewarming hazards \cite{polderman2009,nielsen2013}. Rodent neural-circuit work shows that hibernation-like states can be experimentally induced in a mammal, but no translation to prolonged human torpor follows from that evidence \cite{takahashi2020}.
\section{Survey and Research Gap}
The Survey freezes three evidence boundaries. First, natural hibernation is an integrated physiological program, not passive cooling. Second, human clinical cooling is time-limited and risk-managed. Third, rodent circuit evidence is a mechanistic lead rather than a human efficacy claim. GAP-HT-01 is the absence of a validated human-relevant protocol that jointly controls heat balance, substrate supply, shivering, perfusion, oxygen delivery, electrolytes, coagulation, and rewarming. GAP-HT-02 is the absence of an independently validated safety-and-reversibility framework that separates energy saving from harmful hypothermia.
\section{Idea Selection}
A descriptive comparison of hibernators cannot test human safety. Direct prolonged cooling of healthy volunteers is rejected as unethical. The selected idea is a staged Human Torpor Safety Envelope. It begins with simulation and hardware-in-the-loop fault testing, proceeds only to approved nonhuman translational work, and treats clinically indicated short-duration temperature-management records as observational context. There is no fabricated autonomous idea-search history and no direct escalation to elective human hibernation.
\section{Problem Definition and Hypotheses}
Let $\dot{E}(t)$ be metabolic expenditure, $T_c(t)$ core temperature, and $x(t)$ the vector of perfusion, oxygenation, ECG, electrolyte, acid-base, coagulation, shivering, infection, and neurological-recovery measurements. The relative metabolic change is
\begin{equation}\label{eq:metabolic}
R_E=\frac{\overline{\dot{E}}_{\mathrm{controlled}}-\overline{\dot{E}}_{\mathrm{matched\ normothermia}}}{\overline{\dot{E}}_{\mathrm{matched\ normothermia}}}.
\end{equation}
The proposal does not set $R_E$ to zero and does not interpret a negative value as fuel independence. A state is considered reversible only if all preregistered safety functions remain acceptable during induction, maintenance, and rewarming:
\begin{equation}\label{eq:gate}
G=\indicator\{S_{\mathrm{perfusion}}=S_{\mathrm{oxygen}}=S_{\mathrm{ECG}}=S_{\mathrm{chemistry}}=S_{\mathrm{rewarm}}=1\}.
\end{equation}
Here each $S$ is a locked safety gate defined by domain experts before execution. H1 states that the controller can reduce $R_E$ while retaining $G=1$ in translational testing. H2 states that every sensor-fault or safety breach produces safe abort and rewarming. Failure of either hypothesis rejects escalation.
\section{Study Design and Methods}
Stage 0 develops a physiological digital twin and hardware-in-the-loop controller. It tests sensor latency, heater/cooler fault, perfusion drift, and rewarming failure without exposing living participants. Stage 1 is an approved nonhuman translational protocol with indirect calorimetry, heat-flux measurement, substrate accounting, continuous cardiovascular monitoring, oxygen delivery, blood chemistry, coagulation, shivering, immune markers, and standardized rewarming. Stage 2, if separately approved, is observational analysis of clinically indicated short-duration temperature management; it cannot identify a causal torpor effect.
The control law uses temperature and physiological state, not temperature alone:
\begin{equation}\label{eq:control}
u(t)=\pi\left(T_c(t),x(t),\dot{E}(t),r(t)\right),
\end{equation}
where $u(t)$ is bounded thermal and supportive control, $r(t)$ is the registered reference trajectory, and $\pi$ is frozen before evaluation. Any gate failure overrides $\pi$, terminates induction, and invokes a documented rewarming pathway. No research biopsy, experimental human cooling, or automated clinical action is allowed.
\begin{table}[t]\caption{Staged gates}\centering\footnotesize\begin{tabular}{p{0.23\columnwidth}p{0.67\columnwidth}}\toprule Stage & Requirement \\\midrule 0 & Simulation and fault-injection safety before living-system work.\\ 1 & Nonhuman metabolic reduction plus intact perfusion, chemistry, ECG, coagulation, and rewarming.\\ 2 & Observational, clinically indicated records only; no causal or escalation claim.\\ Release & Independent physiology, critical-care, ethics, and statistics review.\\\bottomrule\end{tabular}\end{table}
\section{Validation and Conditional Outcomes}
Validation locks endpoints, abort thresholds, matched comparators, sensor calibration, and analysis partitions. The primary efficacy measure is $R_E$ from Eq.~\eqref{eq:metabolic}; safety and reversibility are co-primary gates, not secondary caveats. Holdouts include independent runs, operator-blinded controller logs where feasible, and fault-injection scenarios. A successful branch requires metabolic reduction and all gates. A null branch retains normothermia as the preferred condition. A failure branch reports the unsafe boundary and prohibits escalation.
\section{Risks and Human Review}
The largest risk is confusing engineered cooling with evolved hibernation. Additional risks include hypoperfusion, arrhythmia, electrolyte instability, coagulopathy, immune effects, tissue injury, neurological harm, and rewarming shock. The protocol therefore requires expert review at each stage and treats any unresolved risk as a stop condition. Long-duration human torpor remains unproven; even a successful limited study would not establish a fuel-free state or a practical human hibernation technology.
\begin{thebibliography}{00}
\bibitem{geiser2004} F. Geiser, ``Metabolic rate and body temperature reduction during hibernation and daily torpor,'' \emph{Annu. Rev. Physiol.}, vol. 66, pp. 239--274, 2004, doi: 10.1146/annurev.physiol.66.032102.115105.
\bibitem{ruf2015} T. Ruf and F. Geiser, ``Daily torpor and hibernation in birds and mammals,'' \emph{Biol. Rev.}, vol. 90, pp. 891--926, 2015, doi: 10.1111/brv.12137.
\bibitem{polderman2009} K. H. Polderman, ``Mechanisms of action, physiological effects, and complications of hypothermia,'' \emph{Crit. Care Med.}, vol. 37, pp. 1101--1120, 2009, doi: 10.1097/CCM.0b013e318194ba77.
\bibitem{nielsen2013} N. Nielsen \emph{et al.}, ``Targeted temperature management at 33 C versus 36 C after cardiac arrest,'' \emph{N. Engl. J. Med.}, vol. 369, pp. 2197--2206, 2013, doi: 10.1056/NEJMoa1310519.
\bibitem{takahashi2020} T. M. Takahashi \emph{et al.}, ``A discrete neuronal circuit induces a hibernation-like state in rodents,'' \emph{Nature}, vol. 583, pp. 109--114, 2020, doi: 10.1038/s41586-020-2163-6.
\end{thebibliography}
\appendix
\section{Review Checklist}Only frozen Survey sources are used. Before any execution, reviewers must approve species choice, controller limits, endpoint thresholds, nonhuman welfare, clinical-record governance, and the rewarming protocol. This manuscript reports no observed results.
\end{document}
